Much of the data stored in real-world business systems is text, and lossless LLM-based compressors can outperform traditional schemes such as FSST or Brotli by up to $2\times$ on text data. However, the robustness of LLM-based compression to inputs that induce poor model predictions remains understudied. In this paper, we show that LLM prediction failure does not necessarily imply compression failure: unlikely tokens often represent many source bytes, offsetting their higher coding cost. To study this effect, we decompose LLM-based compression into a two-layer system of tokenization and next-token prediction and evaluate an LLM-based compressor on natural, random, and adversarial inputs. A standard prediction attack, \textsc{MinProb}, repeatedly selects the model’s least-probable next token. Although it increases prediction cost by \(5.68\times\) relative to natural text, the resulting output still compresses to \(36.2\%\) of the original size. We therefore introduce a new attack objective, \textsc{MaxSurprisal/Byte}, which targets coding cost per source byte rather than per token. This attack expands the encoded output to \(144.3\%\) of the original size. Our results show that robustness of LLM-based compression must be assessed end-to-end, accounting for both tokenization and prediction.


\begin{comment} % previous version
Much of the data stored in real-world business systems is text.
Large language models (LLMs) have been shown to act as powerful lossless compressors for text data: their next-token probabilities coupled with entropy coding can compress data significantly below its original size.
Such LLM-based compressors can outperform traditional compression schemes such as FSST and Brotli by up to $2\times$ on real-world data.
These benefits come with known drawbacks that have been studied in prior work: lower throughput, higher monetary cost, and the need for specialized hardware.
However, the robustness of these compressors to adversarial inputs remains understudied.
In this paper, we view an LLM-based lossless compressor as a two-layer system consisting of (1) compression through tokenization and (2) compression through entropy coding driven by next-token prediction.
We then evaluate this compressor on either input with random data--model misalignment or challenge the compressor with adversarial inputs that maximize prediction or compression cost with malicious inputs specifically designed to expand the output size.
Crucially, our analysis separates tokenization from LLM prediction and shows that adversarial inputs targeting poor next-token predictions alone, as in the previously proposed \MinProb{} attack, do not necessarily lead to larger output sizes:
even low-probability tokens can represent multiple source bytes and thereby counteract such an attack.
We propose the \MaxSurprisal{} attack, which jointly targets both compression layers. While \MinProb{} increases prediction cost by $5.68\times$ in our experiments, the encoded output still compresses to $36.2\%$ of the original input size. By contrast, our attack increases the encoded size to $144.3\%$ of the original input size.
\end{comment}
